\section{The Laziness Indicators}\label{sec:LazinessIndicators}
The \Java Runtime Library defines several exceptions\index{exception} that are, when get thrown, signs of programmer laziness – or, at the very least, their carelessness. Your code may never throw these exceptions\index{exception} itself, and it must prevent that any method\index{method} called from your code will throw it. Consequently, your code should not contain \lstinline|catch|\index{throw!catch} clauses for any of these exceptions\index{exception}.

These exceptions are:
\begin{itemize}[nosep]
	\item{\lstinline|java.lang.NullPointerException|\index{java.lang!NullPointerException} \autocite{ORACLE_DOC_NULLPOINTEREXCEPTION_CLASS}}
	\item{\lstinline|java.lang.IndexOutOfBoundsException|\index{java.lang!IndexOutOfBoundsException} \autocite{ORACLE_DOC_INDEXOUTOFBOUNDSEXCEPTION_CLASS}
	with the subclasses
	\begin{itemize}[nosep]
		\item{\lstinline|java.lang.ArrayIndexOutOfBoundException|\index{java.lang!ArrayIndexOutOfBoundException} \autocite{ORACLE_DOC_ARRAYINDEXOUTOFBOUNDSEXCEPTION_CLASS}}
		\item{\lstinline|java.lang.StringIndexOutOfBoundException|\index{java.lang!StringIndexOutOfBoundException} \autocite{ORACLE_DOC_STRINGINDEXOUTOFBOUNDSEXCEPTION_CLASS}}
	\end{itemize}}
	\item{\lstinline|java.lang.ArithmethicException|\index{java.lang!ArithmethicException} \autocite{ORACLE_DOC_ARITHMETICEXCEPTION_CLASS}}
%	\item{\lstinline||\index{} \autocite{}}
%	\item{\lstinline||\index{} \autocite{}}
\end{itemize}

Unfortunately the \Java Runtime Library uses these exceptions\index{exception} a lot, mostly to signal invalid arguments.\footnote{However, no one has ever claimed that the authors of the JDK never made questionable decisions.} 

This means you must check the arguments for each call to a method\index{method} that declares to throw one of the listed exceptions. In case of the \lstinline|NullPointerException|\index{java.lang!NullPointerException}, it means also that you need to check each reference before you use it.

Unfortunately, it means also that you cannot use some of the validation methods from the class \lstinline|java.util.Objects|\index{java.util!Objects} \autocite{ORACLE_DOC_OBJECTS_CLASS}:
\begin{itemize}
	\item{All variants of \lstinline|Objects::requireNonNull|\index{java.util!Objects!requireNonNull()} \autocite{ORACLE_DOC_OBJECTS:requireNonNull} throw a \lstinline|NullPointerException|\index{java.lang!NullPointerException} in case the value that is checked is \lstinline|null| – but this is what we want to avoid by checking the value.}
	
	\item{The methods \lstinline|Object::checkIndex|\index{java.util!Objects!checkIndex} \autocite{ORACLE_DOC_OBJECTS:checkIndex}, \lstinline|Object::checkFromToIndex|\index{java.util!Objects!checkFromToIndex()} \autocite{ORACLE_DOC_OBJECTS:checkFromToIndex}, and \lstinline|Object::checkFromIndexSize|\index{java.util!Objects!checkFromIndexSize()} \autocite{ORACLE_DOC_OBJECTS:checkFromIndexSize} throw an \lstinline|IndexOutOfBoundsException|\index{java.lang!IndexOutOfBoundsException} in all their variants  when the value check fails.}
	
%	\item{\lstinline||\index{} \autocite{}}

%	\item{\lstinline||\index{} \autocite{}}
\end{itemize}

You should think about reimplementing these methods to throw a validation exception\index{validation exception} instead. Refer to \autocite{TQUADRAT_ORG_FOUNDATION_OBJECTS}.

And \bdq{Laziness Indicators} because when these exceptions\index{exception} evade to the surface, they will indicate that the programmer was too lazy to implement argument checks for each method\index{method} and each argument.
 
%==============================================================================
% End of file!!
%==============================================================================
